\chapter{旋转、角动量与谱线分裂}
\label{ch:angular}

\section{连续对称性与守恒量}

谐振子每升高一级，振动能量就增加一份 $\hbar\omega$。原子束通过磁场时也出现离散的读数，但这次读数还取决于磁体朝向。
要把沿 $z$ 轴的测量与转过一个角度后的测量联系起来，先要知道怎样转动一个量子态。
空间波函数对无穷小旋转的响应给出轨道角动量；体系的旋转对称性决定相应分量是否守恒。

把一个无自旋粒子的制备绕 $z$ 轴转过小角度 $\epsilon$，探测器的位置固定。
转动后在方位角 $\phi$ 读到的振幅，等于原来在 $\phi - \epsilon$ 处的振幅。
把半径和极角暂时省略，一阶展开为
\begin{equation}
    (\op{U}(\epsilon)\psi)(\phi) = \psi(\phi - \epsilon)
        = \psi(\phi) - \epsilon\,\p_\phi\psi(\phi) + \bigo(\epsilon^2).
\end{equation}
若把这个变化写成 $\op{U}(\epsilon) = \one - \i\epsilon \op{G}/\hbar + \cdots$，
比较 $\epsilon$ 的系数，再用链式法则 $\p_\phi = -y\p_x + x\p_y$，得到绕 $z$ 轴旋转的生成元
\begin{equation}
    \op G = -\i\hbar\,\p_\phi
       = \op X\op P_y - \op Y\op P_x = \op L_z .
\end{equation}
这种确定无穷小变化的算符称为生成元(generator)。
有限维中，幺正性 $\op{U}\adjoint \op{U} = \one$ 的一阶项要求 $\op{G}\adjoint = \op{G}$。
对连续自由度，强连续的单参数幺正群也有自伴生成元（\thmref{thm:qm-stone}）。
一般记 $\epsilon$ 为变换的参数，使乘积 $\epsilon \op{G}$ 具有作用量量纲。

设主动变换 $\op{U}(\epsilon) = \exp(-\i\epsilon\op G/\hbar)$ 保持 Hamilton 算符不变，且 $\op G$ 不显含时间。在光滑波包上对这一不变关系取 $\epsilon = 0$ 处的导数，再用式\eqref{eq:heisenberg}，得到
\begin{equation}
    \op U\adjoint(\epsilon)\op H\op U(\epsilon) = \op H
    \implies [\op G,\op H] = 0
    \implies \dd t\ave{\op G} = 0 .
\end{equation}
这个对称性还使演化保持 $\op{G}$ 的各个谱子空间，因而 $\op{G}$ 的整个测量概率分布守恒。
若生成元本身含时，正确条件是$\p_t\op{G} + (\i/\hbar)[\op{H},\op{G}] = 0$。

在\ssecref{ss:momentum-position}中，平移使每一动量分量积累与位移成正比的相位。写成三维平移算符 $\op{T}(\bm a) = \exp(-\i\bm a \cdot \opbm{P}/\hbar)$，它满足
$\op{T}\adjoint\opbm{X} \op{T} = \opbm{X} + \bm a\,\one$；空间均匀性因此给出动量守恒。
旋转$\op{R}(\bm n,\theta) = \exp(-\i\theta\bm n \cdot \opbm{J}/\hbar)$的生成元满足
\begin{equation}
    [\op{J}_i,\op{J}_j] = \i\hbar\epsilon_{ijk}\op{J}_k.
    \label{eq:Jalgebra}
\end{equation}
无自旋粒子的轨道部分为$\opbm{L} = \opbm{X} \times \opbm{P}$；
有内禀自旋(intrinsic spin)时，旋转同时作用于轨道和自旋，生成元为$\opbm{J} = \opbm{L} + \opbm{S}$。
记 $\opbm{X} = (\op{X},\op{Y},\op{Z})$。由轨道角动量的分量表达式可直接验证式\eqref{eq:Jalgebra}。
写 $\op{L}_x = \op{Y}\op{P}_z - \op{Z}\op{P}_y$、$\op{L}_y = \op{Z}\op{P}_x - \op{X}\op{P}_z$。
展开 $[\op{L}_x,\op{L}_y]$ 的四项，不同坐标与动量彼此对易，只留下
\[
    [\op{L}_x,\op{L}_y] = \op{Y}[\op{P}_z,\op{Z}]\op{P}_x + \op{X}[\op{Z},\op{P}_z]\op{P}_y
        = \i\hbar(\op{X}\op{P}_y - \op{Y}\op{P}_x) = \i\hbar \op{L}_z .
\]
自旋的同一关系已由 Pauli 算符的乘法得到。轨道算符与自旋算符分别作用在空间和内部自由度上，交叉对易子为零，两部分就能直接相加。

质量为 $M$ 的粒子在静标势中的 Hamilton 算符是 $\op H = \opbm P^2/(2M) + V(\opbm X)$，动能与 $\opbm L$ 对易。
由 $[V,\op P_j] = \i\hbar\p_j V$ 和 Ehrenfest 关系可知，角动量均值的变化率等于平均力矩，
\begin{equation}
    \dd{t}\ave{\opbm L} = \ave{\opbm X \times [-\grad V(\opbm X)]}.
\end{equation}
若静势能只依赖距离，力沿径向，每个角动量分量都与 $\op H$ 对易，定态可以按 $\ell,m$ 分类。

\section{角动量谱}

知道角动量产生旋转以后，可以回到束流实验：给定角动量的大小，沿一个选定的轴究竟能分出几束？
确定束数要找出在同一大小下有哪些允许的投影。为此先找能同时有确定值的两个量，再考察怎样从一个投影态走到另一个。

角动量(angular momentum)算符 $\opbm{J}$ 的平方是标量算符 $\opbm{J}^2 = \op{J}_x^2 + \op{J}_y^2 + \op{J}_z^2$。由式\eqref{eq:Jalgebra}，
\begin{equation}
    [\opbm{J}^2,\op{J}_i]
    = \sum_k\bigp{\op{J}_k[\op{J}_k,\op{J}_i] + [\op{J}_k,\op{J}_i]\op{J}_k} = 0,
\end{equation}
因为反对称的$\epsilon_{kil}$与关于$k,l$对称的算符组合缩并为零。
在每个有限维角动量多重态中，由\thmref{thm:qm-compatible}可选用 $\opbm J^2$ 和 $\op J_z$ 的共同本征基。每个基矢都可以用这两个量的确定读数来标记。

\begin{definition}{角动量量子数与升降算符}{qm-angular-labels}
    用无量纲数$\lambda,m$暂记共同本征值：
    \begin{subequations}
        \begin{alignat}{1}
            \opbm{J}^2\bigket{\lambda,m} &= \hbar^2\lambda\bigket{\lambda,m},\\
            \op{J}_z\bigket{\lambda,m} &= \hbar m\bigket{\lambda,m}.
        \end{alignat}
    \end{subequations}
    定义$\op{J}_+ = \op{J}_x + \i \op{J}_y$、$\op{J}_- = \op{J}_x - \i \op{J}_y$，且$\op{J}_+\adjoint = \op{J}_-$。
    若一个非零子空间在三个 $\op{J}_i$ 的作用下封闭，且不含更小的非零共同不变子空间，就称它为不可约(irreducible)角动量多重态。
\end{definition}

\begin{lemma}[升降一步后的投影与范数]
    给定 $\opbm{J}^2,\op{J}_z$ 的归一化共同本征态 $\ket{\lambda,m}$，若 $\op{J}_\pm\ket{\lambda,m}$ 非零，
    它的 $\opbm{J}^2$ 本征值仍为 $\hbar^2\lambda$，投影量子数变为 $m \pm 1$。升降后向量的范数平方为
    $\hbar^2[\lambda - m(m \pm 1)]$。
\end{lemma}
\begin{proof}
由 $\op{J}_\pm$ 的定义和角动量对易关系，
\begin{subequations}
    \begin{alignat}{1}
        [\op{J}_z,\op{J}_\pm] &= \pm\hbar \op{J}_\pm,\\
        \op{J}_\mp \op{J}_\pm &= \opbm{J}^2 - \op{J}_z^2 \mp \hbar \op{J}_z.
    \end{alignat}
\end{subequations}
因 $\opbm{J}^2$ 与各分量对易，升降不改变 $\opbm{J}^2$ 的本征值。由第一式，$\op{J}_z \op{J}_\pm\ket{\lambda,m} = \hbar(m \pm 1)\op{J}_\pm\ket{\lambda,m}$。
在原态上取第二式的期望，左侧就是升降后向量的范数平方，右侧给出所述表达式。
\end{proof}

\begin{theorem}{角动量的大小与投影谱}{qm-angular-spectrum}
    在一个有限维不可约角动量多重态内，有一个整数或半整数 $j \geq 0$，使得
    \begin{subequations}
        \begin{align}
            \opbm{J}^2\ket{j,m} &= \hbar^2j(j + 1)\ket{j,m},\\
            \op{J}_z\ket{j,m} &= \hbar m\ket{j,m},
        \end{align}
    \end{subequations}
    $m = -j,-j + 1,\ldots,j$。
    因此沿一根轴可分辨出 $2j + 1$ 个投影值。
\end{theorem}
\begin{proof}
根据引理，范数平方非负。对固定 $\lambda$，若 $m$ 无限增加或减少，
$\lambda - m(m \pm 1)$最终会变负，因此升降链必须在两端终止。
设最高、最低的$m$分别为$a,b$，端点条件为
\begin{subequations}
    \begin{alignat}{1}
        \op{J}_+\ket{\lambda,a} = 0 &\implies \lambda = a(a + 1),\\
        \op{J}_-\ket{\lambda,b} = 0 &\implies \lambda = b(b - 1).
    \end{alignat}
\end{subequations}
两式相减给出$(a + b)(a - b + 1) = 0$。
由于$a \geq b$，第二个因子为正，故$b = -a$。
记$a = j \geq 0$，便有$\lambda = j(j + 1)$。
从 $-j$ 升到 $j$ 共走 $2j$ 步，每步使 $m$ 增加一，故 $2j$ 必为非负整数。
所得各态的线性组合在三个角动量分量作用下仍留在这条升降链内；不可约性保证这条链已给出整个多重态。
\end{proof}

\begin{corollary}
    由升降后态的范数，并逐态选择相位，可得
    \begin{equation}
        \op{J}_\pm\ket{j,m} = \hbar\sqrt{(j \mp m)(j \pm m + 1)}\ket{j,m \pm 1}.
    \end{equation}
\end{corollary}
普通三维位置自由度的轨道角动量(orbital angular momentum)只有整数$\ell$，
因为标量位置态上的$2\pi$旋转为恒等。单电子轨道用小写正体字母 $\mathrm{s},\mathrm{p},\mathrm{d},\mathrm{f}$ 依次标记 $\ell = 0,1,2,3$。
自旋量子数 $s$ 可以是半整数。
自旋$1/2$的$2\pi$旋转给出$-\one$；若只旋转干涉仪的一条支路，这个相对符号可以被读出。

在图\ref{fig:angular-projection}所示的 $\opbm{L}^2,\op{L}_z$ 共同本征态中，横向分量仍有量子涨落。
\begin{figure}[htbp]
    \centering\tikzfigure{angular-projection}
    \caption{$\ell = 2$ 时球半径为 $\sqrt{6}\hbar$，五条纬圈的高度对应五个 $\op{L}_z$ 本征值。点线标出 $m = 1$ 态的轴向投影，其读数为 $m\hbar = \hbar$。}\label{fig:angular-projection}
\end{figure}

\subsection{轨道角动量的方向波函数}
\label{ss:qm-spherical-harmonics}
若测量粒子的位置方向，轨道态 $\ket{\ell,m}$ 在立体角 $\d\Omega = \sin\theta\d\theta\d\phi$
内的概率为 $\abs{Y_{\ell m}(\theta,\phi)}^2\d\Omega$。方向振幅 $Y_{\ell m}$ 称为球谐函数(spherical harmonic)。
在极角 $\theta$、方位角 $\phi$ 的球坐标中，$\bm r = r\uvec r$ 与梯度叉乘消去了梯度的径向分量，故
\begin{equation}
    \opbm L = -\i\hbar\,\bm r \times \grad
       = -\i\hbar\biggp{\uvec\phi\,\p_\theta
            - \frac{\uvec\theta}{\sin\theta}\p_\phi}.
\end{equation}
取它在 $z$ 轴的分量得 $\op L_z = -\i\hbar\p_\phi$，取 $x,y$ 分量的组合则有
\begin{equation}
    \op L_\pm = \op L_x \pm \i\op L_y
        = \hbar\e^{\pm\i\phi}\bigp{\pm\p_\theta + \i\cot\theta\,\p_\phi}.
\end{equation}
因而 $\op L_\pm$ 使方位因子 $\e^{\i m\phi}$ 的 $m$ 增减一。
将这两个微分算符代入 $\opbm L^2 = \op L_+\op L_- + \op L_z^2 - \hbar\op L_z$，得
\begin{equation}
    \opbm L^2 = -\hbar^2\biggb{\frac{1}{\sin\theta}\p_\theta\bigp{\sin\theta\,\p_\theta}
                       + \frac{1}{\sin^2\theta}\p_\phi^2}.
    \label{eq:qm-angular-laplacian}
\end{equation}
由 $\op L_zY_{\ell m} = \hbar mY_{\ell m}$ 可写 $Y_{\ell m} = f(\cos\theta)\e^{\i m\phi}$。
再令 $x = \cos\theta$，$\opbm L^2$ 的本征方程成为
\begin{equation}
    \dd{x}\biggp{(1 - x^2)\dv{f}{x}}
        + \biggp{\ell(\ell + 1) - \frac{m^2}{1 - x^2}}f = 0 .
\end{equation}
在两极正则的解由缔合 Legendre 函数给出。取 Condon--Shortley 相位约定，对 $0 \le m \le \ell$ 写成
\begin{subequations}
    \begin{align}
        Y_{\ell m} &= \sqrt{\frac{2\ell + 1}{4\pi}\frac{(\ell - m)!}{(\ell + m)!}}\,
            P_\ell^m(\cos\theta)\e^{\i m\phi},\\
        P_\ell^m(x) &= (-1)^m(1 - x^2)^{m/2}\dv[m]{P_\ell(x)}{x}.
    \end{align}
\end{subequations}
其中的 Legendre 多项式可以由 Rodrigues 公式计算，
\begin{equation}
    P_\ell(x) = \frac{1}{2^\ell\ell!}\dd[\ell]{x}(x^2 - 1)^\ell .
\end{equation}
负投影的方向振幅满足 $Y_{\ell,-m} = (-1)^m Y_{\ell m}\conj$。
沿整个方位角积分得到 $2\pi\delta_{mm'}$；固定 $m$ 后，极角积分给出不同 $\ell$ 的正交性。两者合起来有
\begin{equation}
    \int_{S^2}Y_{\ell m}\conj Y_{\ell' m'}\d\Omega = \delta_{\ell\ell'}\delta_{mm'}.
\end{equation}

\begin{example}{轨道角动量的联合测量}{qm-cartesian-angular-readout}
    一个无自旋粒子的归一化波函数为 $\psi(\bm r) = x^2 f(r)$，其中 $r = \abs{\bm r}$。同时测量 $\opbm L^2$ 和 $\op L_z$ 会得到哪几对读数，各有多大概率？若只保留测得 $L_z = 0$ 的粒子，再测得 $L^2 = 6\hbar^2$ 的概率是多少？
    \par\smallskip\textbf{解}\quad
    把 $r^2$ 并入径向因子，角向因子便是 $(x/r)^2 = \sin^2\theta\cos^2\phi$。把 $\cos^2\phi$ 拆成 $(1 + \cos 2\phi)/2$，只会出现 $m = 0,\pm2$。其中与 $\phi$ 无关的部分还要按 $\ell$ 分开。用 $P_2(u) = (3u^2 - 1)/2$ 可得
    \[
        \sin^2\theta\cos^2\phi
        = \frac{1 - P_2(\cos\theta)}{3} + \frac12\sin^2\theta\cos 2\phi .
    \]
    第一项含常数和 $\ell = 2$ 的 Legendre 多项式。最后一项可以拆成带 $\e^{\pm2\i\phi}$ 的两个球谐。换成上文的归一化球谐并将系数的模平方和归一，得到方向振幅
    \[
        \frac{\sqrt{5}}{3}Y_{00} - \frac13Y_{20}
        + \frac{Y_{22} + Y_{2,-2}}{\sqrt{6}} .
    \]
    因此 $(L^2,L_z) = (0,0)$ 的概率为 $5/9$，$(6\hbar^2,0)$ 的概率为 $1/9$，$(6\hbar^2,\pm2\hbar)$ 的两个结果各占 $1/6$。
    单独筛选 $L_z = 0$ 时，两种不同 $\ell$ 的正交分量都通过，总概率为 $5/9 + 1/9 = 2/3$。在通过的粒子中再测得 $L^2 = 6\hbar^2$ 的概率为 $(1/9)/(2/3) = 1/6$。
\end{example}

\section{角动量合成}
单个角动量的大小和投影都有了标签。原子同时有轨道角动量和自旋，转动整个原子时两者一起转动。
分开测量两个投影，与测量总角动量的大小和投影，给出两种不同的态标签。
两个自旋均为 $1/2$ 的粒子给出了一个例子：不同乘积态可具有相同的总投影，测量总角动量的大小还能区分它们的某些叠加态。

设有两个可区分的自旋，各自的 $z$ 分析器给出 $\up$ 或 $\dn$。
按第一个、第二个自旋的次序记录，两台分析器可以给出 $(\up,\up),(\up,\dn),(\dn,\up),(\dn,\dn)$ 四种结果。
在\dfnref{dfn:qm-tensor-product}的张量积空间中，每一对结果对应一个乘积基矢。例如 $\ket{\up,\dn} = \ket{\up}_1 \otimes \ket{\dn}_2$。
若两边独立制备在 $\ket{a}$ 和 $\ket{b}$，
联合振幅是两边振幅的乘积，内积也相应满足
$\braket{a',b'}{a,b} = \braket{a'}{a}\braket{b'}{b}$。
在联合空间里，$\op{J}_{1z}$ 简写 $\op{J}_{1z} \otimes \one_2$，$\op{J}_{2z}$ 简写 $\one_1 \otimes \op{J}_{2z}$。
分别作用于两个因子的算符可以交换作用次序，因此它们对易。
例如 $\ket{\up,\dn}$ 给出的总投影为 $\hbar/2 - \hbar/2 = 0$。
另一个态 $\ket{\dn,\up}$ 也给出零，总投影相同的结果还可以按总角动量的大小区分。
\begin{theorem}{两个角动量的合成}{qm-angular-addition}
    固定两个子系统的角动量量子数 $j_1,j_2$，取各自一个不可约多重态。
    在其张量积空间中，总角动量 $\opbm{J} = \opbm{J}_1 + \opbm{J}_2$ 的量子数依次为
    \begin{equation}
        j = \abs{j_1 - j_2},\abs{j_1 - j_2} + 1,\ldots,j_1 + j_2 .
    \end{equation}
    每个这样的 $j$ 出现一个多重态，每个多重态中 $m = m_1 + m_2$ 从 $-j$ 取到 $j$。
\end{theorem}
\begin{proof}
    设 $j_1 \geq j_2$，对允许的非负投影 $m$，数出有多少乘积态满足 $m_1 + m_2 = m$，记为 $d_m$。
    当 $m \leq j_1 - j_2$ 时，$m_2$ 从 $-j_2$ 到 $j_2$ 均可取，故 $d_m = 2j_2 + 1$。
    当 $j_1 - j_2 \leq m \leq j_1 + j_2$ 时，为使 $m_1 = m - m_2 \leq j_1$，
    $m_2$ 只能从 $m - j_1$ 取到 $j_2$，于是 $d_m = j_1 + j_2 - m + 1$。
    超过最高投影后 $d_m = 0$。

    从最高投影态开始逐次下降，取走一整个多重态；其正交补仍在各 $\op{J}_i$ 作用下封闭，可继续作同样的分解。
    每个 $j \geq m + 1$ 的多重态各给 $d_m$ 和 $d_{m+1}$ 贡献一个态，只有 $j = m$ 的多重态仅给前者贡献一个。
    所以 $j = m$ 的出现次数是 $d_m - d_{m+1}$。代入各投影态的数目，这个差在 $j_1 - j_2 \leq m \leq j_1 + j_2$ 时为一，在更小的 $m$ 处为零。
\end{proof}

\begin{definition}{Clebsch--Gordan系数}{qm-cg}
    固定 $j_1,j_2$ 后，分别测量两个角动量的 $z$ 投影，得到非耦合基(uncoupled basis)，
    \begin{equation}\label{eq:qm-uncoupled-angular-momentum-basis}
        \ket{m_1,m_2} = \ket{j_1,m_1}_1 \otimes \ket{j_2,m_2}_2 .
    \end{equation}
    改测总角动量的平方 $\opbm J^2$ 和总投影 $\op J_z$，得到耦合基(coupled basis)。
    基矢写作 $\ket{j_1,j_2;j,m}$，分号前标出参加合成的两个角动量，分号后给出总量子数。

    耦合基矢在乘积基上的分量称为 Clebsch--Gordan 系数，
    \begin{equation}
        C^{jm}_{j_1m_1,j_2m_2}
        = \bigbraket{m_1,m_2}{j_1,j_2;j,m}.
    \end{equation}
    只有$m = m_1 + m_2$的系数可能非零。
    同时列出两组基时，用 $\ket{j,m}_J$ 简记耦合态，角标 $J$ 表明两个标签都描述总角动量。
\end{definition}

\begin{example}{两个自旋的耦合基}{qm-two-spins}
    两个可区分粒子的自旋均为 $1/2$。求总角动量平方和总投影的共同本征态，用两个自旋的 $z$ 投影乘积基表示。
    \par\smallskip\textbf{解}\quad
    对两个自旋$1/2$，总磁量子数$m = 1$只有一个态$\ket{\up,\up}$，
    故它是$j = 1$的最高态。
    用$\op{J}_- = \op{S}_{1-} + \op{S}_{2-}$作用，
    \[
        \op{J}_-\ket{\up,\up}
        = \hbar\bigp{\ket{\dn,\up} + \ket{\up,\dn}}
        = \hbar\sqrt2\ket{1,0}_J.
    \]
    再下降一次得到$\ket{1,-1}_J = \ket{\dn,\dn}$。
    $m = 0$子空间中与$\ket{1,0}_J$正交的归一化态
    $(\ket{\up,\dn} - \ket{\dn,\up})/\sqrt2$
    被$\op{J}_+$和$\op{J}_-$同时湮灭，因此$j = 0$。
\end{example}

把这四个耦合基矢用乘积基展开，得到
\begin{subequations}
    \begin{alignat}{1}
        \ket{1,1}_J &= \ket{\up,\up},\qquad \ket{1,-1}_J = \ket{\dn,\dn},\\
        \ket{1,0}_J &= \frac{\ket{\up,\dn} + \ket{\dn,\up}}{\sqrt2},\qquad
            \ket{0,0}_J = \frac{\ket{\up,\dn} - \ket{\dn,\up}}{\sqrt2}.
    \end{alignat}
\end{subequations}
前三个是三重态(triplet)，最后一个是单态(singlet)。例如 $\ket{1,0}_J$ 的两个数给出确定的 $j = 1,m = 0$。若分别测两个自旋的 $z$ 投影，$(\up,\dn)$ 和 $(\dn,\up)$ 各有一半概率。
附录\ref{app:cg}把这四个态的系数整理成表，并列出轨道角动量与自旋 $1/2$ 耦合时的系数。
借助$\opbm{J}_1 \cdot \opbm{J}_2 = (\opbm{J}^2 - \opbm{J}_1^2 - \opbm{J}_2^2)/2$，
交换模型$\op{H} = (A/\hbar^2)\opbm{S}_1 \cdot \opbm{S}_2$的三重态能量为$A/4$、单态为$-3A/4$。这里 $A$ 是有能量量纲的耦合常数。
因此交换耦合把四维自旋空间分成三重态与单态两个能量子空间。

\section{Zeeman效应}
\label{sec:zeeman}
轨道和自旋合成后，给定 $j$ 的各个 $m_j$ 在零场中可以同能。外加磁场选定了一根轴，
沿轴的磁矩投影不同，原子与场的相互作用能也不同。于是原来的一条谱线能够分裂成几条。
要预言光谱仪上各条线的位置，必须把发光前、后两个能级各自的磁能变化都算进去。

\subsection{磁场中的光谱测量}
\label{ss:zeeman-experiment}

Zeeman实验研究外磁场是否改变发光原子的谱线。
把光源放在磁极之间，用高分辨率光谱仪比较加场与不加场时的谱线，
并通过偏振分析器分别观察不同偏振成分。
实验不仅测量原有谱线是否展宽或分裂，还测量各成分随磁场的频移及其偏振。
这些量共同约束原子磁矩与角动量的关系\cite{zeemanlecture}。

光谱仪直接读到的是上下能级之差。
用下标 $\hi$ 和 $\lo$ 分别标记上、下能级。若它们的零场能量为 $E_\hi^{(0)}$ 和 $E_\lo^{(0)}$，则无场频率为
$\nu_0 = (E_\hi^{(0)} - E_\lo^{(0)})/h$。
加场后某一允许跃迁的频率为
\begin{equation}
    \nu_{m_\hi,m_\lo}(B)
    = \nu_0 + \frac{\Delta E_\hi(m_\hi,B) - \Delta E_\lo(m_\lo,B)}h.
    \label{eq:zeeman-line}
\end{equation}
求谱线位置，先分别求上、下能级的磁子态能量，再按允许跃迁配对取能量差。

\subsection{磁矩与弱场能移}
\label{ss:zeeman-weak}

固定原子核近似下，电子磁矩及线性磁相互作用为
\begin{subequations}
    \begin{align}
        \opbm{\mu} &= -\frac{\muB}{\hbar}(\opbm{L} + g_{\mathrm{s}}\opbm{S}),\\
        \op{H}_{\mathrm{Z}} &= \frac{\muB B}{\hbar}(\op{L}_z + g_{\mathrm{s}}\op{S}_z).
    \end{align}
\end{subequations}
这里$g_{\mathrm{s}} \simeq 2$为最低阶电子自旋因子，精密实验还需反常磁矩修正。
设零场能级已被自旋轨道作用按$j$分开，且Zeeman能远小于相邻$j$能级间隔。
在量子数 $\ell,s,j$ 固定且 $j > 0$ 的子空间，旋转对称性(rotational symmetry)要求$\opbm{\mu}$的投影正比于$\opbm{J}$。

令 $\op\Pi_j$ 投影到固定 $\ell,s,j$ 的多重态，用比例系数 $g_J$ 表示上述投影关系，
\begin{equation}
    \op\Pi_j(\opbm L + g_{\mathrm{s}}\opbm S)\op\Pi_j = g_J\op\Pi_j\opbm J\op\Pi_j .
\end{equation}
因为$\opbm{L} + g_{\mathrm{s}}\opbm{S} = \opbm{J} + (g_{\mathrm{s}} - 1)\opbm{S}$，
点乘$\opbm{J}$并在$\ket{j,m_j}$中取期望，得到
\[
    g_J\hbar^2j(j + 1)
    = \hbar^2j(j + 1)
    + (g_{\mathrm{s}} - 1)\frac{\hbar^2}{2}
    \bigb{j(j + 1) + s(s + 1) - \ell(\ell + 1)}.
\]

该投影关系也可从角动量合成看出：
在固定$\ell,s,j$的一个不可约子空间内，矢量算符的三个分量
与$\op{J}_x,\op{J}_y,\op{J}_z$具有相同的旋转变换关系，差别只剩一个与$m_j$无关的系数。
将$\opbm{S} \cdot \opbm{J} = (\opbm{J}^2 + \opbm{S}^2 - \opbm{L}^2)/2$代入后，
除以$\hbar^2j(j + 1)$便得到
\begin{subequations}
    \begin{align}
        g_J &= 1 + (g_{\mathrm{s}} - 1)\frac{j(j + 1) + s(s + 1) - \ell(\ell + 1)}{2j(j + 1)},\\
        \Delta E &= \muB B g_Jm_j.
    \end{align}
\end{subequations}

这里$g_J$称为Landé因子，$m_j$是$\op{J}_z/\hbar$的本征值。
上述 $g_J$ 公式适用于 $j > 0$。非简并的 $j = 0$ 态由旋转对称性给出零磁矩，一阶能移为零。
\subsection{谱线分裂与偏振}
\label{ss:zeeman-lines}

将上下能级的线性能移代入式\eqref{eq:zeeman-line}，得
\begin{equation}
    \nu_{m_\hi,m_\lo}
    = \nu_0 + \frac{\muB B}{h}
        \bigp{g_\hi m_\hi - g_\lo m_\lo}.
\end{equation}
在电偶极近似中，允许的磁量子数差为$m_\hi - m_\lo = 0,\pm1$。
差为零的跃迁称为$\pi$成分，差为$\pm1$的称为$\sigma$成分。
沿磁场方向观察时，$\pi$偶极沿观察方向不辐射，
两个$\sigma$成分对应相反的圆偏振；
垂直磁场观察时，$\pi$与$\sigma$成分呈相互正交的线偏振。
偏振分析由此提供了独立于频移的跃迁标签。

对于$s = 0$，$g_\hi = g_\lo = 1$，于是
\begin{equation}
    \nu_{m_\hi,m_\lo} - \nu_0
    = \frac{\muB B}{h}(m_\hi - m_\lo)
    \in \BigB{0,\ \frac{\muB B}{h},\  -\frac{\muB B}{h}}.
\end{equation}
多组不同$m_\hi,m_\lo$可能具有同一个频移，
所以即使原子有多个磁子态，正常Zeeman效应仍形成三个频率组（图\ref{fig:zeeman-triplet}）。
各组总强度则应把落在该组内的跃迁按初态布居加权求和。
\begin{figure}[htbp]
    \centering\tikzfigure{zeeman-triplet}
    \caption{正常 Zeeman 三重线：从 $j_\hi = 1$ 的三个磁子态向同一 $j_\lo = 0$ 态发射。横向分列三个通道，箭头的高度对应光子能量；图中放大了磁分裂。}\label{fig:zeeman-triplet}
\end{figure}

\begin{example}{双重态谱线的 Zeeman 分裂}{qm-zeeman-doublet}
    一个闭壳层外有一个价电子的原子处于弱磁场 $B$ 中。对于电偶极跃迁 $\,{}^2\mathrm{P}_{1/2} \to {}^2\mathrm{S}_{1/2}$，忽略核自旋并取 $g_{\mathrm{s}} = 2$。求各个 $\pi$ 和 $\sigma$ 成分相对于零场谱线的频移，以及不同频率的数目。
    \par\smallskip\textbf{解}\quad
    在光谱项上标出总角动量后，能级记为 ${}^{2S+1}L_J$。中间按总轨道角动量量子数 $L = 0,1,2,3$ 依次填入大写正体字母 $\mathrm{S},\mathrm{P},\mathrm{D},\mathrm{F}$。
    左上角是由总自旋量子数 $S$ 决定的多重度 $2S + 1$，右下角是总角动量量子数 $J$。
    闭壳层不贡献角动量，故 $L = \ell$、$S = s = 1/2$、$J = j$。
    代入Landé公式，得到
    \begin{subequations}
        \begin{alignat}{2}
            g_\hi &= 1 + \frac{3/4 + 3/4 - 2}{2(3/4)}
                = \frac23,
            &\quad m_\hi &= \pm\frac12,\\
            g_\lo &= 1 + \frac{3/4 + 3/4}{2(3/4)}
                = 2,
            &\quad m_\lo &= \pm\frac12.
        \end{alignat}
    \end{subequations}
    两个$\pi$跃迁满足$m_\hi = m_\lo$，频移为$\mp2\muB B/(3h)$；
    两个$\sigma$跃迁连接相反的$m$，频移为$\pm4\muB B/(3h)$。
    四种允许跃迁对应四个不同的上下能级差，形成四条谱线。
    四个频率的强度还取决于磁子态布居、偏振和角动量耦合系数。
\end{example}

若零场相邻$j$多重态的间距为$\Delta_\mathrm{fs}$，
上述投影计算要求磁混合矩阵元远小于这个间距。
磁场增加到两者可比时，$\opbm{J}^2$不再近似守恒，
应在共同的固定$m_j$子空间内对角化自旋轨道项与Zeeman项的和。

\subsectionstar{Paschen--Back极限}

当 Zeeman 能量超过自旋轨道劈裂，而正比于 $B^2$ 的抗磁能移仍可忽略时，轨道与自旋可分别按 $m_\ell,m_s$ 标记。在这个 Paschen--Back 极限，乘积态 $\ket{\ell,m_\ell} \otimes \ket{s,m_s}$ 的磁能移为
\begin{equation}
    \Delta E \simeq \muB B(m_\ell + g_{\mathrm{s}} m_s).
\end{equation}

\section{宇称与选择定则}
列出磁场中全部能级后，若把上、下能级任意配对，会预言出实验中看不到的线。
一次电偶极跃迁要有非零的偶极矩阵元。把原子作空间反演或绕量子化轴旋转，便能在计算径向积分之前，
先判断哪些配对的矩阵元必定为零。

\begin{definition}{空间反演与宇称}{qm-parity}
    宇称算符 $\op{\mathcal P}$ 实现空间反演，将相对原点的位置反号，对标量位置波函数有 $(\op{\mathcal P}\psi)(\bm x) = \psi(-\bm x)$。
    $\op{\mathcal P}$ 幺正且 $\op{\mathcal P}^2 = \one$，本征值为 $+1$ 或 $-1$，分别称为偶宇称和奇宇称(parity)。
\end{definition}
空间反演使动量反号，即 $\op{\mathcal P}\opbm P\op{\mathcal P}^{-1} = -\opbm P$。位置 $\opbm X$ 也反号，故轨道角动量 $\opbm{L} = \opbm{X} \times \opbm{P}$ 保持不变。
若$[\op{H},\op{\mathcal P}] = 0$，非简并能量本征态必有确定宇称；简并时可在该子空间内选共同本征基。
一般制备态可以叠加两种宇称的分量。普通中心势轨道态的宇称为$(-1)^\ell$。

把原子绕 $z$ 轴旋转，电偶极矩的两个横向分量 $\op d_x,\op d_y$ 彼此混合，纵向分量 $\op d_z$ 保持不变。
无穷小旋转给出 $[\op J_z,\op d_x] = \i\hbar\op d_y$ 和 $[\op J_z,\op d_y] = -\i\hbar\op d_x$。
因此对组合 $\op d_x \pm \i\op d_y$ 取同一对易子，只会把该组合乘以 $\pm\hbar$。据此定义电偶极矩的球分量，
\begin{subequations}\label{eq:dipole-spherical-components}
    \begin{gather}
        \op d_0 = \op d_z,\qquad
        \op d_{\pm1} = \mp\frac{\op d_x \pm \i\op d_y}{\sqrt2},\\
        {}[\op J_z,\op d_q] = \hbar q\op d_q .
    \end{gather}
\end{subequations}
绕 $z$ 轴旋转时，每个新分量只乘一个相位。
这样的三个分量构成秩为 1 的球张量。这里的秩指一组算符在旋转下怎样变换。

\begin{definition}{球张量算符}{qm-spherical-tensor}
    一组$2k + 1$个算符$\op{T}_q^{(k)}$（$q = -k,\ldots,k$）
    若满足
    \begin{subequations}
        \begin{alignat}{1}
            [\op{J}_z,\op{T}_q^{(k)}] &= \hbar q\op{T}_q^{(k)},\\
            [\op{J}_\pm,\op{T}_q^{(k)}] &= \hbar\sqrt{(k \mp q)(k \pm q + 1)}\,\op{T}_{q\pm1}^{(k)},
        \end{alignat}
    \end{subequations}
    就称为秩$k$的不可约球张量(irreducible spherical tensor)。
    $q$ 标记各球分量，第二式给出旋转生成元如何联系相邻的分量。
\end{definition}

只用第一条对易式就能筛去一些谱线。设
$M = \mel{j',m'}{\op{T}_q^{(k)}}{j,m}$，将 $\op{J}_z$ 分别作用在左矢和右矢上，得到
\[
    \mel{j',m'}{[\op{J}_z,\op{T}_q^{(k)}]}{j,m}
        = \hbar(m' - m)M = \hbar qM.
\]
所以只有 $m' = m + q$ 时 $M$ 才可能非零。例如沿 $z$ 的电场只耦合偶极的
$\op{d}_z = \op{d}_0$ 分量，吸收这样的光保持 $m$。其余磁量子数的跃迁可用改变偏振来探测。

\begin{theorem}{Wigner--Eckart定理}{qm-wigner-eckart}
    对秩$k$的不可约球张量，其角动量矩阵元可以分解为
    \begin{equation}
        \mel{\alpha',j',m'}{\op{T}_q^{(k)}}{\alpha,j,m}
        = \frac{C^{j'm'}_{jm,kq}}{\sqrt{2j' + 1}}
        \langle\alpha',j'\|\op{T}^{(k)}\|\alpha,j\rangle.
    \end{equation}
    $\alpha,\alpha'$表示除角动量之外的其余态标签，如主量子数。
    双竖线表示约化矩阵元(reduced matrix element)，它与$m,m',q$无关，归一化由上式约定。
\end{theorem}

这里 $C^{j'm'}_{jm,kq}$ 是将 $j$ 和 $k$ 合成 $j'$ 的系数，采用\dfnref{dfn:qm-cg}的约定。
定理将旋转决定的角依赖与动力学决定的约化矩阵元分开；
其一般证明可参照Littlejohn的讲义，比较公式时需统一约化矩阵元的归一化约定\cite{littlejohnwe}。
三角规则(triangle rule) $\abs{j - k} \leq j' \leq j + k$和$m' = m + q$来自角动量合成。
对电偶极$\opbm{d}$，$k = 1$且宇称为奇，因此必须改变宇称，
$\Delta j = 0,\pm1$但$0 \leftrightarrow 0$禁止，$\Delta m = 0,\pm1$。
角动量选择定则筛出允许通道；允许通道的跃迁强度还取决于径向积分，它也可能为零。

宇称选择规则也可直接检查。若
$\op{\mathcal P}\ket{i} = \pi_i\ket{i}$、$\op{\mathcal P}\ket{f} = \pi_f\ket{f}$，其中$\pi_i,\pi_f = \pm1$，
且$\op{\mathcal P} \op{d}_q\op{\mathcal P}^{-1} = -\op{d}_q$，则
\[
    -\mel{f}{\op{d}_q}{i}
    = \mel{f}{\op{\mathcal P}\op{d}_q\op{\mathcal P}^{-1}}{i}
    = \pi_f\pi_i\mel{f}{\op{d}_q}{i}.
\]
同宇称时矩阵元等于自身的相反数，故矩阵元为零；
相反宇称时该约束自动满足。

\sectionstar{时间反演与Kramers简并}
空间反演能迫使某些跃迁的强度为零。回到零磁场，还可以问另一个光谱问题：
即使原子周围的环境破坏了旋转对称性，半整数自旋的能级为何仍可能成对出现？
将一次运动反向播放时，角动量变号。若反向运动服从同一个 Hamilton 算符，就能用它寻找同能的伙伴。

\begin{definition}{反幺正变换与时间反演}{qm-time-reversal}
    反幺正(antiunitary)变换 $\op{\Theta}$ 是共轭线性的双射，满足
    $\braket{\op{\Theta} u}{\op{\Theta} v} = \braket{v}{u}$。
    时间反演(time reversal)由这样的算符表示，保持位置不变，使动量和角动量反号。
\end{definition}
自旋$1/2$可取$\op{\Theta} = -\i\op{\sigma}_y \op{K}$，其中$\op{K}$是在$z$基的复共轭，故$\op{\Theta}^2 = -\one$。
\begin{theorem}{Kramers简并}{qm-kramers}
    若自伴Hamilton算符$\op{H}$在反幺正变换$\op{\Theta}$下不变，
    即$\op{\Theta} \op{H}\op{\Theta}^{-1} = \op{H}$，且$\op{\Theta}^2 = -\one$，
    则每个离散能量本征态$\ket{n}$都有与它正交的同能伙伴$\op{\Theta}\ket{n}$。
\end{theorem}
\begin{proof}
    在反幺正内积关系中取$\ket{u} = \ket{n}$、$\ket{v} = \op{\Theta}\ket{n}$，得到
    \[
        \braket{\op{\Theta} n}{\op{\Theta}^2n} = \braket{\op{\Theta} n}{n}.
    \]
    左侧又等于$-\braket{\op{\Theta} n}{n}$，所以这个内积为零。
    若$\op{H}\ket{n} = E_n\ket{n}$，由$E_n$为实数与$[\op{H},\op{\Theta}] = 0$，
    \begin{equation}
        \op{H}\op{\Theta}\ket{n} = \op{\Theta} \op{H}\ket{n} = E_n\op{\Theta}\ket{n}.
    \end{equation}
    这两个正交态由此构成同一能量下的一对态。
\end{proof}

整个实验作时间反演时，外磁场也随之反向；固定外磁场通常破坏时间反演对称性。

\begin{example}{旋转对称体系的基态}{qm-rotational-ground-multiplet}
    某体系的 Hamilton 算符在任意空间旋转下不变。若基态有确定的总角动量量子数 $j$，这个对称性对其简并度有何要求？仅凭旋转对称性能否确定 $j$？
    \par\smallskip\textbf{解}\quad
    同一 $j$ 多重态内的 $m = -j,\ldots,j$ 能量相同，故基态至少有 $2j + 1$ 个线性无关的态。哪个 $j$ 多重态能量最低，取决于具体的 Hamilton 算符、粒子统计和边界条件。
\end{example}
\reading{角动量的升降与合成核对MIT 8.05讲义9、10及Feynman第III卷第18章\cite{mitangular,feynman18}。
    球坐标下的梯度与 Laplacian 可查《电动力学》卷附录“曲线坐标系”，球谐函数的相位与归一化见\href{https://dlmf.nist.gov/14.30}{NIST 数学函数手册第14.30节}。
    下一章用这些代数解释氢谱中更大的简并。}
